
\documentclass{article}
\usepackage{colortbl}
\usepackage{makecell}
\usepackage{multirow}
\usepackage{supertabular}

\begin{document}

\newcounter{utterance}

\centering \large Interaction Transcript for game `cladder', experiment `full\_v1.5\_default', episode 7000 with score{-}Qwen3.5{-}9B{-}sp.
\vspace{24pt}

{ \footnotesize  \setcounter{utterance}{1}
\setlength{\tabcolsep}{0pt}
\begin{supertabular}{c@{$\;$}|p{.15\linewidth}@{}p{.15\linewidth}p{.15\linewidth}p{.15\linewidth}p{.15\linewidth}p{.15\linewidth}}
    \# & \multicolumn{2}{c}{Player} && \multicolumn{2}{c}{Game Master} \\
    \hline

    \theutterance \stepcounter{utterance}  
    & & & \multicolumn{4}{p{0.6\linewidth}}{
        \cellcolor[rgb]{0.9,0.9,0.9}{
            \makecell[{{p{\linewidth}}}]{
                \texttt{\tiny{[P1$\langle$GM]}}
                \texttt{You are an expert in causal inference. The following question is not a typical commonsense query, but rather a meticulously designed question created by a professor specializing in causal inference, intended to assess the students' mastery of the course content.} \\
\texttt{Imagine a self{-}contained, hypothetical world with only the following conditions, and without any unmentioned factors or causal relationships: Jyka has a direct effect on gyzp and hwax. Hwax has a direct effect on lirg. Gyzp has a direct effect on lirg. For those who are not jyka, the probability of lirg is 63\%. For those who are jyka, the probability of lirg is 49\%. For those who are jyka, would it be more likely to see lirg if the individual was not jyka?} \\
\texttt{Start your answer with 'Yes' or 'No', followed by additional reasoning or evidence to support your explanation.} \\
            }
        }
    }
     \\ \\

    \theutterance \stepcounter{utterance}  
    & \multicolumn{4}{p{0.6\linewidth}}{
        \cellcolor[rgb]{0.561,0.737,0.561}{
            \makecell[{{p{\linewidth}}}]{
                \texttt{\tiny{[P1$\rangle$GM]}}
                \texttt{No} \\
\\ 
\texttt{**Reasoning:**} \\
\\ 
\texttt{To answer this question, we must carefully parse the conditional probabilities provided in the hypothetical world and compare the likelihood of the outcome `lirg` under two specific counterfactual scenarios for an individual who is `jyka`.} \\
\\ 
\texttt{1.  **Identify the Baseline Probabilities:**} \\
\texttt{    The problem explicitly states the probability of `lirg` based on the status of `jyka`:} \\
\texttt{    *   \$P(\textbackslash{}text\{lirg\} \textbackslash{}mid \textbackslash{}text\{jyka\} = \textbackslash{}text\{true\}) = 49\textbackslash{}\%\$} \\
\texttt{    *   \$P(\textbackslash{}text\{lirg\} \textbackslash{}mid \textbackslash{}text\{jyka\} = \textbackslash{}text\{false\}) = 63\textbackslash{}\%\$} \\
\\ 
\texttt{2.  **Analyze the Question's Premise:**} \\
\texttt{    The question asks: "For those who are **jyka**, would it be more likely to see **lirg** if the individual was **not jyka**?"} \\
\texttt{    *   **Subject:** An individual who is currently `jyka` (\$\textbackslash{}text\{jyka\} = \textbackslash{}text\{true\}\$).} \\
\texttt{    *   **Counterfactual Condition:** The individual is hypothetically changed to `not jyka` (\$\textbackslash{}text\{jyka\} = \textbackslash{}text\{false\}\$).} \\
\texttt{    *   **Target Outcome:** The probability of `lirg`.} \\
\\ 
\texttt{3.  **Compare the Probabilities:**} \\
\texttt{    We need to compare the probability} \\
            }
        }
    }
    & & \\ \\

    \theutterance \stepcounter{utterance}  
    & & & \multicolumn{2}{p{0.3\linewidth}}{
        \cellcolor[rgb]{0.9,0.9,0.9}{
            \makecell[{{p{\linewidth}}}]{
                \texttt{\tiny{[GM$|$GM]}}
                \texttt{No} \\
            }
        }
    }
    & & \\ \\

    \theutterance \stepcounter{utterance}  
    & & & \multicolumn{2}{p{0.3\linewidth}}{
        \cellcolor[rgb]{0.9,0.9,0.9}{
            \makecell[{{p{\linewidth}}}]{
                \texttt{\tiny{[GM$|$GM]}}
                \texttt{yes} \\
            }
        }
    }
    & & \\ \\

    \theutterance \stepcounter{utterance}  
    & & & \multicolumn{2}{p{0.3\linewidth}}{
        \cellcolor[rgb]{0.9,0.9,0.9}{
            \makecell[{{p{\linewidth}}}]{
                \texttt{\tiny{[GM$|$GM]}}
                \texttt{game\_result = LOSE} \\
            }
        }
    }
    & & \\ \\

\end{supertabular}
}

\end{document}
